% =====================================================================
% sections_eval_results.tex  --  Updated Evaluation Methodology and
% Benchmark Results sections using the comprehensive n=30 data.
% 
% Data sources:
%   finance / gpt-5.4-mini: 30 pooled (runs 20260720, 20260728)
%   finance / gpt-5.4:      30/arm (run 20260729T080229)
%   agenticpay / gpt-5.4-mini: 10/arm (run 20260727)
%   agenticpay / gpt-5.4:      10/arm (run 20260728)
%
% To integrate: replace sec:method and sec:results in main.tex,
% or \input this file at the appropriate location.
% =====================================================================

\section{Evaluation Methodology}\label{sec:method}

\subsection{Conformance and goal achievement}\label{sec:conform}

\textbf{Protocol conformance (Set~A).}
Let $T$ be a run's event trace and $M_r$ the projected EFSM for role~$r$.
A message event $e = (s, r, \ell, v)$ \emph{conforms} iff it is accepted by
$M_s$ (structurally legal send) and satisfies any associated refinement
predicate $\phi_\ell(v)$ from the guard sidecar.
The \emph{violation rate} of a trace is the fraction of events rejected by
at least one role's monitor.
Conformance is deterministic, per-message, and requires no LLM judge.

\textbf{Goal achievement (Set~B).}
Each case defines $k$ goals $G_1, \ldots, G_k$, each anchored to a specific
message edge $(s, r, \ell)$ and a payload predicate $\phi_i(v)$.
A trial \emph{succeeds} iff all $k$ goals are satisfied in at least one
of its allowed attempts.
Goals are graded under two rungs:
\emph{strict} requires exact $(s, r, \ell)$ match plus $\phi_i$;
\emph{role-pair} requires $(s, r)$ match plus $\phi_i$ under any label.
Configurations shown the protocol vocabulary are graded strictly;
the bare configuration (never shown the labels) is graded on the
role-pair rung---a fairness correction adopted after audit
(\S\ref{sec:fairaudit}).

\textbf{Design principle.}
The separation into Set~A (structural conformance) and Set~B (task
completion) is not incidental---it is the empirical instantiation of
the theoretical claim that \emph{type safety and progress are distinct
guarantees}.
A system can achieve zero violations while failing its goals (the
protocol is followed but the session runs out of steps), or achieve
all goals while accumulating violations (correct outcome, incorrect
path).
Reporting both sets is necessary to distinguish safety from liveness.

\subsection{Consequence-graded violations: the severity scale S0--S4}\label{sec:severity}

Each surviving deviation is graded by consequence on five levels:
\textbf{S0}~benign (reordering of independent messages; not counted);
\textbf{S1}~waste (duplicate message, costs tokens, breaks nothing);
\textbf{S2}~broken obligation (skipped read or wrong ordering);
\textbf{S3}~non-termination (session exhausts its budget);
\textbf{S4}~unauthorized irreversible act (filing before approval).
The grading is computed against the partial order of data and authority
dependencies declared in the protocol.
Empirical validation: across all configurations and runs, every attempt
containing an S2+ violation went on to fail its goal.

\subsection{Experimental configurations}\label{sec:configs}

Each experiment runs the same task under five \emph{configurations}
(arms), forming a ladder of increasing protocol enforcement.
All share the same intent, role descriptions, model, and retry budget;
only the protocol information and enforcement mechanism differ.

\begin{table}[h]
\centering\small
\caption{The five experimental configurations. Each adds one mechanism
to the previous. ``Protocol info'' = what is pasted into the agent's prompt;
``Enforcement'' = what the runtime does with off-contract messages.}
\label{tab:configs}
\begin{tabular}{@{}llll@{}}
\toprule
\textbf{Arm} & \textbf{Protocol info} & \textbf{Enforcement} & \textbf{Scheduling} \\
\midrule
\texttt{bare} & Intent only (natural language) & None & Round-robin \\
\texttt{maf} & Full validated global type as text & None (MAF GroupChat) & LLM speaker-select \\
\texttt{min\_llmvalid} & Projected local type per role & None (observe only) & Round-robin \\
\texttt{gate} & Projected local type per role & Gate rejects violations & Round-robin \\
\texttt{sched} & Projected local type per role & Gate rejects violations & EFSM-driven \\
\bottomrule
\end{tabular}
\end{table}

The ladder isolates three orthogonal effects:
(i)~\emph{protocol knowledge} (bare$\to$min\_llmvalid): does telling the
agent the protocol help?
(ii)~\emph{enforcement} (min\_llmvalid$\to$gate): does blocking wrong
messages help beyond telling?
(iii)~\emph{scheduling} (gate$\to$sched): does polling only enabled roles
help beyond blocking?
The MAF arm tests an alternative delivery mechanism (global text +
LLM-based speaker selection) that represents the dominant industry
pattern.

\subsection{Statistical methodology}

All confidence intervals are \textbf{Wilson score intervals} (95\%,
$z = 1.96$), appropriate for binomial proportions and well-behaved at
boundary values (0\% and 100\%).
Token counts are the primary efficiency metric; wall-clock times are
reported but labeled indicative due to parallel execution and rate-limit
contention.
Each case uses balanced branch allocation (equal trials per protocol
branch) to prevent branch-composition confounds.

% =====================================================================

\section{Experiments and Results}\label{sec:results}

We evaluate \stjp{} on two multi-agent coordination tasks spanning
different protocol complexities, each run under two model capability
tiers.
The design produces a $2 \times 2 \times 5$ matrix:
2~cases $\times$ 2~models $\times$ 5~arms, with $n \geq 26$ per cell
on the primary case.

\subsection{Cases}

\textbf{Finance} (6~roles, branching choice, 6~goals).
A quarterly revenue report produced by Fetcher, RevenueAnalyst,
ExpenseAnalyst, Writer, TaxVerifier, and TaxSpecialist.
Revenue above \$50k triggers an audit branch requiring TaxSpecialist
analysis and TaxVerifier approval before the report can be written.
The protocol exercises fan-out notifications, value-dependent branching,
and convergence of parallel analysis streams.

\textbf{AgenticPay Settlement} (4~roles, linear, 4~goals).
An escrow-mediated trade between Buyer, Seller, Escrow, and Carrier.
The protocol is linear (no branching) but requires strict sequencing:
deposit $\to$ ship $\to$ release $\to$ settle.
A circular-wait variant (each role waits for the other) demonstrates
the deadlock the session type prevents.

\subsection{Models}

\textbf{GPT-5.4 Mini}: a cost-efficient model where enforcement effects
are most visible (weaker models make more coordination errors).
\textbf{GPT-5.4}: a frontier-capability model that tests whether
enforcement remains necessary as model quality improves.
Both are deployed on Azure AI Foundry.

\subsection{Main results: finance ($n{=}30$)}\label{sec:finance-results}

\begin{table}[t]
\centering\small
\caption{Finance case: goal completion (strict), violations, and cost.
$n = 30$ per arm on GPT-5.4; $n = 26$--$30$ pooled on GPT-5.4 Mini
(two independent runs).
Wilson 95\% CIs in brackets.
Token savings computed relative to bare.
Bold = best in column.}
\label{tab:finance}
\begin{tabular}{@{}l cc cc cc@{}}
\toprule
& \multicolumn{2}{c}{\textbf{Goal completion (\%)}}
& \multicolumn{2}{c}{\textbf{Violation rate (\%)}}
& \multicolumn{2}{c}{\textbf{Tokens/trial}} \\
\cmidrule(lr){2-3}\cmidrule(lr){4-5}\cmidrule(lr){6-7}
\textbf{Arm} & \textbf{Mini} & \textbf{Full}
             & \textbf{Mini} & \textbf{Full}
             & \textbf{Mini} & \textbf{Full} \\
\midrule
\texttt{bare}
  & 42.1\rlap{$^{\dagger}$} {\scriptsize [23,64]}
  & 0.0\rlap{$^{\dagger}$} {\scriptsize [0,11]}
  & 97.3 & 80.1
  & 135{,}879 & 58{,}007 \\
\texttt{maf}
  & 6.7 {\scriptsize [1.8,21]}
  & \textbf{100.0} {\scriptsize [89,100]}
  & 28.7 & \textbf{0.0}
  & 31{,}519 & 25{,}083 \\
\texttt{min\_llmvalid}
  & 90.0 {\scriptsize [74,97]}
  & 50.0 {\scriptsize [32,68]}
  & 21.9 & \textbf{0.0}
  & 50{,}200 & 87{,}613 \\
\texttt{gate}
  & 52.0 {\scriptsize [34,70]}
  & \textbf{100.0} {\scriptsize [89,100]}
  & \textbf{0.0} & \textbf{0.0}
  & 54{,}633 & 37{,}915 \\
\texttt{sched}
  & \textbf{100.0} {\scriptsize [89,100]}
  & \textbf{100.0} {\scriptsize [89,100]}
  & \textbf{0.0} & \textbf{0.0}
  & \textbf{12{,}652} & \textbf{12{,}729} \\
\bottomrule
\end{tabular}
\\[4pt]
{\footnotesize $\dagger$ Role-pair rung (no protocol vocabulary); strict is N/A.}
\end{table}

Five findings emerge from Table~\ref{tab:finance}.

\textbf{(1)~SCHED is model-invariant.}
The EFSM-scheduled arm achieves 100\% goal completion on both models
at near-identical cost (${\sim}$12.7k tokens/trial, 9.5~calls/trial).
No other arm achieves this: every alternative's success rate is a
function of model capability.
SCHED's calls/trial (11~high branch, 8~standard) matches the
protocol's theoretical minimum message count, confirming the scheduler
polls \emph{only} roles with enabled SEND transitions.

\textbf{(2)~Knowledge-only regresses on the stronger model.}
\texttt{min\_llmvalid} drops from 90\% (Mini) to 50\% (Full)---a
40-percentage-point regression ($p < 0.01$, Fisher exact).
The mechanism: GPT-5.4 produces zero protocol violations (agents
follow message ordering correctly) but its verbose outputs consume
more turns before reaching the terminal goal (G6 = 50\%).
This demonstrates that \emph{type safety (conformance) and progress
(task completion) are orthogonal guarantees}: the protocol is followed,
but the session runs out of steps.
Token cost rises 75\% ($50{,}200 \to 87{,}613$), confirming verbosity
as the mechanism.

\textbf{(3)~Gate and MAF scale positively with model capability.}
\texttt{gate}: $52\% \to 100\%$ (+48pp).
\texttt{maf}: $6.7\% \to 100\%$ (+93pp).
The stronger model's improved instruction-following makes enforcement
unnecessary \emph{for correctness}---but not for cost (Finding~4).

\textbf{(4)~SCHED is 4.6--6.9$\times$ cheaper than every alternative at 100\%.}
Even when all enforced arms succeed, cost differs dramatically:
\texttt{sched} at 12.7k tokens vs.\
\texttt{maf} at 25.1k ($2.0\times$),
\texttt{gate} at 37.9k ($3.0\times$),
\texttt{bare} at 58.0k ($4.6\times$),
\texttt{min\_llmvalid} at 87.6k ($6.9\times$).
The mechanism is call count: sched makes 9.5 calls/trial (protocol
minimum) vs.\ gate's 34.0 (round-robin polling of all 6 roles even
when 5 are blocked).

\textbf{(5)~The enforcement gradient is non-monotonic across models.}
On Mini: $\text{bare}(42\%) < \text{maf}(7\%) < \text{gate}(52\%) < \text{min\_llmvalid}(90\%) < \text{sched}(100\%)$.
On Full: $\text{bare}(0\%) < \text{min\_llmvalid}(50\%) < \{\text{maf}, \text{gate}, \text{sched}\}(100\%)$.
Knowledge-only inverts position between models; bare itself
\emph{regresses} ($42\% \to 0\%$ role-pair); only SCHED is
monotonically at the top.

\subsection{Per-goal analysis: why min\_llmvalid regresses}\label{sec:pergol}

\begin{table}[t]
\centering\small
\caption{Per-goal strict completion (\%), finance, GPT-5.4.
\texttt{min\_llmvalid} passes all structural goals (G1--G3, G5) but
fails G4 (strict label mismatch; role-pair G4 = 100\%) and G6
(session budget exhaustion).
All other protocol-aware arms achieve 100\% on every goal.}
\label{tab:pergol}
\begin{tabular}{@{}lcccccc@{}}
\toprule
\textbf{Arm} & \textbf{G1} & \textbf{G2} & \textbf{G3} & \textbf{G4} & \textbf{G5} & \textbf{G6} \\
\midrule
\texttt{bare} (rp) & 50 & 53 & 0 & 83 & 100 & 100 \\
\texttt{maf} & 100 & 100 & 100 & 100 & 100 & 100 \\
\texttt{min\_llmvalid} & 100 & 100 & 100 & 50 & 100 & 50 \\
\texttt{gate} & 100 & 100 & 100 & 100 & 100 & 100 \\
\texttt{sched} & 100 & 100 & 100 & 100 & 100 & 100 \\
\bottomrule
\end{tabular}
\end{table}

The regression mechanism is visible in Table~\ref{tab:pergol}.
\texttt{min\_llmvalid} achieves perfect structural conformance (zero
violations) but fails G6 (final report) in 50\% of trials.
Trace analysis reveals the pattern: without scheduling, the runtime
polls all 6~roles every round via round-robin; GPT-5.4's longer
outputs (each valid, each protocol-conforming) consume the step budget
before Writer produces the terminal \texttt{GenerateReport} message.
On Mini, terser outputs complete within budget (90\%); on Full,
verbosity exhausts it.
This is not a model defect but a \emph{structural limitation of
observation-only enforcement}: a monitor can verify that every message
is correct, but it cannot force progress.
The scheduler resolves this by eliminating idle polls entirely (9.5
calls vs.\ 79.0), leaving ample budget for the terminal goal.

\subsection{AgenticPay settlement ($n{=}10$)}\label{sec:agenticpay}

\begin{table}[t]
\centering\small
\caption{AgenticPay settlement (4~roles, linear protocol, $n = 10$
per arm).
Wilson 95\% CIs at $n{=}10$ are wide ([72,100] at 100\%); the
differentiation is on Mini.}
\label{tab:agenticpay}
\begin{tabular}{@{}l cc cc@{}}
\toprule
& \multicolumn{2}{c}{\textbf{Strict goal completion (\%)}}
& \multicolumn{2}{c}{\textbf{Tokens/trial}} \\
\cmidrule(lr){2-3}\cmidrule(lr){4-5}
\textbf{Arm} & \textbf{Mini} & \textbf{Full}
             & \textbf{Mini} & \textbf{Full} \\
\midrule
\texttt{bare}
  & 0.0 {\scriptsize [0,28]} & 0.0 {\scriptsize [0,28]}
  & 98{,}260 & 82{,}544 \\
\texttt{maf}
  & 0.0\rlap{$^*$} {\scriptsize [0,28]} & 100.0 {\scriptsize [72,100]}
  & 49{,}307 & 22{,}965 \\
\texttt{min\_llmvalid}
  & 55.6 {\scriptsize [27,81]} & 100.0 {\scriptsize [72,100]}
  & 97{,}552 & 46{,}408 \\
\texttt{gate}
  & 100.0 {\scriptsize [72,100]} & 100.0 {\scriptsize [72,100]}
  & 102{,}630 & 48{,}497 \\
\texttt{sched}
  & \textbf{100.0} {\scriptsize [72,100]} & \textbf{100.0} {\scriptsize [72,100]}
  & \textbf{44{,}022} & \textbf{15{,}441} \\
\bottomrule
\end{tabular}
\\[4pt]
{\footnotesize $^*$ 100\% role-pair (correct participants), 0\% strict (wrong labels).}
\end{table}

On the simpler linear protocol, GPT-5.4 achieves 100\% across all
protocol-aware arms (Table~\ref{tab:agenticpay})---confirming that a
four-role linear sequence is within a frontier model's coordination
capacity even without enforcement.
The differentiation appears on Mini, where the gradient separates:
\texttt{bare} (0\%) $<$ \texttt{maf} (0\% strict, 100\% role-pair)
$<$ \texttt{min\_llmvalid} (55.6\%)
$<$ \texttt{gate}/\texttt{sched} (100\%).
Cost follows the same pattern: \texttt{sched} is 81\% cheaper than bare
on Full (15.4k vs.\ 82.5k) and 55\% cheaper on Mini.

\textbf{Deadlock demonstration.}
A circular-wait variant of the protocol (\texttt{unchecked\_skills}:
skill-based agents without shared protocol) produces \emph{zero events}
across all trials---agents wait indefinitely for messages that never
arrive.
The Scribble checker rejects the circular-wait global type before any
agent runs, at zero token cost.

\subsection{Model-scaling analysis}\label{sec:scaling}

\begin{table}[t]
\centering\small
\caption{Model-scaling effects on finance (primary case, $n \geq 26$).
$\Delta$ = change from Mini to Full.
Only SCHED is invariant on both axes (correctness and cost).}
\label{tab:scaling}
\begin{tabular}{@{}l rr c rr c@{}}
\toprule
& \multicolumn{2}{c}{\textbf{Goal compl.\ (\%)}} & &
  \multicolumn{2}{c}{\textbf{Tokens/trial}} & \\
\cmidrule(lr){2-3}\cmidrule(lr){5-6}
\textbf{Arm} & \textbf{Mini} & \textbf{Full} & $\Delta$
             & \textbf{Mini} & \textbf{Full} & $\Delta$ \\
\midrule
\texttt{bare}        & 42.1 &   0   & $-42$pp    & 135{,}879 & 58{,}007 & $-57\%$ \\
\texttt{maf}         &  6.7 & 100.0 & $+93$pp    &  31{,}519 & 25{,}083 & $-20\%$ \\
\texttt{min\_llmvalid} & 90.0 &  50.0 & $\mathbf{-40}$pp & 50{,}200 & 87{,}613 & $+75\%$ \\
\texttt{gate}        & 52.0 & 100.0 & $+48$pp    &  54{,}633 & 37{,}915 & $-31\%$ \\
\texttt{sched}       & 100.0 & 100.0 & $\mathbf{0}$pp & 12{,}652 & 12{,}729 & $+0.6\%$ \\
\bottomrule
\end{tabular}
\end{table}

Table~\ref{tab:scaling} reveals a striking non-monotonicity.
Most arms improve with model capability (gate: $+48$pp, maf: $+93$pp),
consistent with the intuition that stronger models need less
enforcement.
But \texttt{min\_llmvalid} \emph{regresses} ($-40$pp, $p < 0.01$)---the
\textbf{capability-without-efficiency trap}---and even bare regresses
($42\% \to 0\%$ role-pair) as the stronger model's assertiveness
breaks the accidental serialization that rescued the weaker one.
The explanation is that a stronger model with no enforcement produces
longer, more detailed outputs that are individually correct (zero
violations) but collectively exhaust the session budget.
This result falsifies the hypothesis that enforcement becomes
unnecessary as models scale, and establishes that SCHED---the only
arm combining type safety (gate) with progress (scheduling)---is
the only configuration providing model-independent guarantees.

\textbf{Three regimes of model scaling:}
\begin{enumerate}\itemsep0pt
\item \emph{Enforcement-independent} (SCHED): 100\%/100\%, constant cost.
      The guarantee is structural and does not depend on model quality.
\item \emph{Positively scaling} (gate, maf): improve with capability but
      require the frontier model. Unsuitable for deployments where model
      quality cannot be guaranteed.
\item \emph{Negatively scaling} (min\_llmvalid): degrades with capability.
      Knowledge alone is a trap: the model follows the protocol perfectly
      while exhausting the budget.
\end{enumerate}

\subsection{Cost analysis: the scheduling dividend}\label{sec:cost}

\begin{table}[t]
\centering\small
\caption{Cost decomposition, finance, GPT-5.4 ($n = 30$).
SCHED achieves optimal calls by polling only enabled roles;
its calls/trial matches the protocol's minimum message count.}
\label{tab:cost}
\begin{tabular}{@{}l rrr r@{}}
\toprule
\textbf{Arm} & \textbf{Calls/trial} & \textbf{Tok/call} & \textbf{Tok/trial} & \textbf{Savings} \\
\midrule
\texttt{bare}           & 45.2 & 1{,}283 & 58{,}007 & --- \\
\texttt{maf}            & 13.7 & 1{,}831 & 25{,}083 & 56.8\% \\
\texttt{min\_llmvalid}  & 79.0 & 1{,}109 & 87{,}613 & $-51\%$ \\
\texttt{gate}           & 34.0 & 1{,}115 & 37{,}915 & 34.6\% \\
\texttt{sched}          &  9.5 & 1{,}340 & 12{,}729 & \textbf{78.1\%} \\
\bottomrule
\end{tabular}
\end{table}

Two mechanisms produce the cost differences (Table~\ref{tab:cost}).
\textbf{Call reduction}: SCHED makes 9.5 calls/trial (the protocol's
minimum: 11 messages on the high branch, 8 on standard, averaged) vs.\
gate's 34.0 (round-robin polls all 6~roles every round) and bare's
45.2 (retries on failure inflate the count).
Tokens-per-call is roughly constant across arms (${\sim}$1.1--1.8k),
so the entire cost difference comes from \emph{making fewer calls, not
shorter ones}.
\textbf{Budget exhaustion}: \texttt{min\_llmvalid}'s 79 calls (vs.\ 9.5)
are all valid---each produces a conforming message---but most are wasted
on roles that cannot advance the protocol, consuming budget that would
otherwise reach G6.

The per-branch breakdown confirms optimality:
\texttt{sched} uses 11.0 calls on high (protocol has 11 messages) and
8.0 on standard (protocol has 8).
\texttt{gate} uses 40.0 and 28.0 respectively---the difference is pure
idle polling overhead.

\subsection{Cross-case comparison}\label{sec:crosscase}

\begin{table}[t]
\centering\small
\caption{Cross-case comparison (both cases, GPT-5.4 Mini).
The scheduler's cost advantage grows with protocol complexity:
from 55\% savings (4-role linear) to 91\% (6-role branching).}
\label{tab:crosscase}
\begin{tabular}{@{}l cc cc@{}}
\toprule
& \multicolumn{2}{c}{\textbf{AgenticPay (4 roles, linear)}}
& \multicolumn{2}{c}{\textbf{Finance (6 roles, branching)}} \\
\cmidrule(lr){2-3}\cmidrule(lr){4-5}
\textbf{Metric} & \texttt{gate} & \texttt{sched}
               & \texttt{gate} & \texttt{sched} \\
\midrule
Goal completion & 100\% & 100\% & 52\% & \textbf{100\%} \\
Violations & 0 & 0 & 0 & 0 \\
Tokens/trial & 102{,}630 & 44{,}022 & 54{,}633 & \textbf{12{,}652} \\
Savings vs.\ bare & 0\% & 55\% & 60\% & \textbf{91\%} \\
Calls/trial & 84.2 & 36.0 & 54.9 & \textbf{9.5} \\
\bottomrule
\end{tabular}
\end{table}

Table~\ref{tab:crosscase} shows the scheduling dividend \emph{grows}
with protocol complexity.
On the linear 4-role protocol, both gate and sched achieve 100\%
(the linear structure has no out-of-turn ambiguity), but sched saves
55\% on tokens by eliminating idle polls.
On the 6-role branching protocol, gate degrades to 52\% (round-robin
polling exhausts retry budget on roles blocked by the branch) while
sched maintains 100\% and saves 91\%.
The mechanism is clear: with more roles, round-robin polling wastes
proportionally more calls on blocked roles; the EFSM scheduler scales
near-optimally because it polls only the (at most 2) enabled roles
at each state.

\subsection{Discussion: implications for multi-agent system design}

The results support three design principles:

\textbf{Principle~1: Type safety does not imply progress.}
\texttt{min\_llmvalid} on GPT-5.4 demonstrates this formally:
zero protocol violations, 50\% goal completion.
A session type guarantees that every message that \emph{happens} is
correct; it does not guarantee that the \emph{right} messages happen
in time.
Progress requires a scheduler that restricts polling to enabled roles,
preventing budget waste on idle ones.

\textbf{Principle~2: Model scaling is not a substitute for structural
enforcement.}
The non-monotonic regression of \texttt{min\_llmvalid} (90\%$\to$50\%)
shows that ``just use a bigger model'' can make coordination
\emph{worse} when the system lacks enforcement.
A stronger model that is more verbose, more confident, and produces
longer outputs will exhaust finite resources faster when nothing
constrains it.

\textbf{Principle~3: The cheapest correct system is the most constrained.}
SCHED's 78--91\% cost savings come not from compressing outputs but
from \emph{not making calls whose answers cannot advance the protocol}.
The protocol itself is the scheduling oracle: the EFSM's enabled-set
at each state is exactly the set of roles worth polling.
